% Options for packages loaded elsewhere
% Options for packages loaded elsewhere
\PassOptionsToPackage{unicode}{hyperref}
\PassOptionsToPackage{hyphens}{url}
%
\documentclass[
  ignorenonframetext,
  aspectratio=169,
  russian,
]{beamer}
\newif\ifbibliography
\usepackage{pgfpages}
\setbeamertemplate{caption}[numbered]
\setbeamertemplate{caption label separator}{: }
\setbeamercolor{caption name}{fg=normal text.fg}
\beamertemplatenavigationsymbolshorizontal
% Prevent slide breaks in the middle of a paragraph
\widowpenalties 1 10000
\raggedbottom
\AtBeginPart{
  \frame{\partpage}
}
\AtBeginSection{
  \ifbibliography
  \else
    \frame{\sectionpage}
  \fi
}
\AtBeginSubsection{
  \frame{\subsectionpage}
}
\usepackage{iftex}
\ifPDFTeX
  \usepackage[T1]{fontenc}
  \usepackage[utf8]{inputenc}
  \usepackage{textcomp} % provide euro and other symbols
\else % if luatex or xetex
  \usepackage{unicode-math} % this also loads fontspec
  \defaultfontfeatures{Scale=MatchLowercase}
  \defaultfontfeatures[\rmfamily]{Ligatures=TeX,Scale=1}
\fi
\usepackage{lmodern}

\ifPDFTeX\else
  % xetex/luatex font selection
\fi
% Use upquote if available, for straight quotes in verbatim environments
\IfFileExists{upquote.sty}{\usepackage{upquote}}{}
\IfFileExists{microtype.sty}{% use microtype if available
  \usepackage[]{microtype}
  \UseMicrotypeSet[protrusion]{basicmath} % disable protrusion for tt fonts
}{}


\usepackage{longtable,booktabs,array}
\usepackage{calc} % for calculating minipage widths
\usepackage{caption}
% Make caption package work with longtable
\makeatletter
\def\fnum@table{\tablename~\thetable}
\makeatother
\usepackage{graphicx}
\makeatletter
\newsavebox\pandoc@box
\newcommand*\pandocbounded[1]{% scales image to fit in text height/width
  \sbox\pandoc@box{#1}%
  \Gscale@div\@tempa{\textheight}{\dimexpr\ht\pandoc@box+\dp\pandoc@box\relax}%
  \Gscale@div\@tempb{\linewidth}{\wd\pandoc@box}%
  \ifdim\@tempb\p@<\@tempa\p@\let\@tempa\@tempb\fi% select the smaller of both
  \ifdim\@tempa\p@<\p@\scalebox{\@tempa}{\usebox\pandoc@box}%
  \else\usebox{\pandoc@box}%
  \fi%
}
% Set default figure placement to htbp
\def\fps@figure{htbp}
\makeatother



\ifLuaTeX
\usepackage[bidi=basic,provide=*]{babel}
\else
\usepackage[bidi=default,provide=*]{babel}
\fi
% get rid of language-specific shorthands (see #6817):
\let\LanguageShortHands\languageshorthands
\def\languageshorthands#1{}


\setlength{\emergencystretch}{3em} % prevent overfull lines

\providecommand{\tightlist}{%
  \setlength{\itemsep}{0pt}\setlength{\parskip}{0pt}}



 

\usepackage[]{csquotes}

\IfFileExists{plex-otf.sty}{
  %% Full TeXlive
  \usepackage[
  math,
    RM={Scale=0.94},SS={Scale=0.94},SScon={Scale=0.94},TT={Scale=MatchLowercase,FakeStretch=0.9},DefaultFeatures={Ligatures=Common}
   ]{plex-otf}
}{
  %% TinyTeX
  \usepackage{libertine}
}

%%% Load theme
% https://deic.uab.cat/~iblanes/beamer_gallery/
\IfFileExists{beamerthemegotham.sty}{
  %% Full TeXlive
  \usetheme{gotham}
  \gothamset{
    numbering=totalpagenumber,
    parttocframe default=off,
    sectiontocframe default=off,
    subsectiontocframe default=off,
  }
}{
  %% TinyTeX
  \usetheme{Madrid}
}
\makeatletter
\@ifpackageloaded{caption}{}{\usepackage{caption}}
\AtBeginDocument{%
\ifdefined\contentsname
  \renewcommand*\contentsname{Содержание}
\else
  \newcommand\contentsname{Содержание}
\fi
\ifdefined\listfigurename
  \renewcommand*\listfigurename{Список иллюстраций}
\else
  \newcommand\listfigurename{Список иллюстраций}
\fi
\ifdefined\listtablename
  \renewcommand*\listtablename{Список таблиц}
\else
  \newcommand\listtablename{Список таблиц}
\fi
\ifdefined\figurename
  \renewcommand*\figurename{Рисунок}
\else
  \newcommand\figurename{Рисунок}
\fi
\ifdefined\tablename
  \renewcommand*\tablename{Таблица}
\else
  \newcommand\tablename{Таблица}
\fi
}
\@ifpackageloaded{float}{}{\usepackage{float}}
\floatstyle{ruled}
\@ifundefined{c@chapter}{\newfloat{codelisting}{h}{lop}}{\newfloat{codelisting}{h}{lop}[chapter]}
\floatname{codelisting}{Список}
\newcommand*\listoflistings{\listof{codelisting}{Листинги}}
\makeatother
\makeatletter
\makeatother
\makeatletter
\@ifpackageloaded{caption}{}{\usepackage{caption}}
\@ifpackageloaded{subcaption}{}{\usepackage{subcaption}}
\makeatother

\usepackage{bookmark}
\IfFileExists{xurl.sty}{\usepackage{xurl}}{} % add URL line breaks if available
\urlstyle{same}
\hypersetup{
  pdftitle={Моделирование неравновесной агрегации},
  pdfauthor={Жукова Арина Александровна; Садова Диана Алексеевна; Агаев Арсений Валерьевич; Диденко Дмитрий Владимирович},
  pdflang={ru-RU},
  hidelinks,
  pdfcreator={LaTeX via pandoc}}


\title{Моделирование неравновесной агрегации}
\subtitle{Этап 2. Алгоритмы и численные методы}
\author{Жукова Арина Александровна \and Садова Диана
Алексеевна \and Агаев Арсений Валерьевич \and Диденко Дмитрий
Владимирович}
\date{2026-04-11}

\begin{document}
\frame{\titlepage}

\renewcommand*\contentsname{Содержание}
\begin{frame}[allowframebreaks]
  \frametitle{Содержание}
  \setcounter{tocdepth}{2}
  \tableofcontents
\end{frame}
\setcounter{tocdepth}{2}
\tableofcontents
}

\section{1. Информация}\label{ux438ux43dux444ux43eux440ux43cux430ux446ux438ux44f}

\begin{frame}{1.1 Докладчики}
\phantomsection\label{ux434ux43eux43aux43bux430ux434ux447ux438ux43aux438}
\textbf{Жукова Арина Александровна} студент 3 курса, 1132239120@rudn.ru

\textbf{Садова Диана Алексеевна} студент 3 курса, 1132239118@rudn.ru

\textbf{Агаев Арсений Валерьевич} студент 3 курса, 1032221668@rudn.ru

\textbf{Диденко Дмитрий Владимирович} студент 3 курса,
1032230532@rudn.ru

\emph{Российский университет дружбы народов} \emph{Кафедра прикладной
информатики и теории вероятностей}
\end{frame}

\section{2. Вводная
часть}\label{ux432ux432ux43eux434ux43dux430ux44f-ux447ux430ux441ux442ux44c}

\begin{frame}{2.1 Актуальность}
\phantomsection\label{ux430ux43aux442ux443ux430ux43bux44cux43dux43eux441ux442ux44c}
\begin{itemize}
\tightlist
\item
  Процессы неравновесной агрегации встречаются повсеместно: образование
  сажи, рост кристаллов, электрический пробой
\item
  Модель DLA (Diffusion Limited Aggregation) --- классический пример
  формирования фрактальных структур
\item
  Для компьютерного моделирования необходимы эффективные алгоритмы
\item
  Правильный выбор алгоритма определяет скорость и точность расчетов
\end{itemize}
\end{frame}

\begin{frame}{2.2 Цели и задачи этапа}
\phantomsection\label{ux446ux435ux43bux438-ux438-ux437ux430ux434ux430ux447ux438-ux44dux442ux430ux43fux430}
\textbf{Цель:} разработать алгоритмическую базу для моделирования DLA и
расчета фрактальной размерности

\textbf{Задачи:} - Описать базовый алгоритм DLA на сетке - Разработать
методы расчета фрактальной размерности - Предложить модификации для
различных физических условий - Обосновать вычислительные оптимизации
\end{frame}

\section{3. Математический
аппарат}\label{ux43cux430ux442ux435ux43cux430ux442ux438ux447ux435ux441ux43aux438ux439-ux430ux43fux43fux430ux440ux430ux442}

\begin{frame}{3.1 Случайные блуждания}
\phantomsection\label{ux441ux43bux443ux447ux430ux439ux43dux44bux435-ux431ux43bux443ux436ux434ux430ux43dux438ux44f}
\textbf{На сетке (четыре направления):}

Вероятность каждого направления равна \(1/4\).

\textbf{В непрерывном пространстве:}

Смещение на шаге \(h\):
\[x_{new} = x + h \cos\alpha, \quad y_{new} = y + h \sin\alpha\] где
\(\alpha\) --- случайный угол от \(0\) до \(2\pi\).
\end{frame}

\begin{frame}{3.2 Фрактальная размерность}
\phantomsection\label{ux444ux440ux430ux43aux442ux430ux43bux44cux43dux430ux44f-ux440ux430ux437ux43cux435ux440ux43dux43eux441ux442ux44c}
\textbf{Метод радиуса гирации:}

\[N \sim R_g^D\]

\[R_g = \sqrt{\frac{1}{N}\sum_{i=1}^N \left[(x_i-x_c)^2 + (y_i-y_c)^2\right]}\]

\textbf{Метод ящиков (box counting):}

\[N(\varepsilon) \sim \varepsilon^{-D}\]

Логарифмическая зависимость: \[\ln N = D \cdot \ln R_g + C\]
\end{frame}

\section{4. Базовый алгоритм
DLA}\label{ux431ux430ux437ux43eux432ux44bux439-ux430ux43bux433ux43eux440ux438ux442ux43c-dla}

\begin{frame}{4.1 Алгоритм роста кластера}
\phantomsection\label{ux430ux43bux433ux43eux440ux438ux442ux43c-ux440ux43eux441ux442ux430-ux43aux43bux430ux441ux442ux435ux440ux430}
\textbf{Шаг 1. Инициализация:} - Создать пустую сетку размера
\(L \times L\) - Поместить затравочную частицу в центр - Счетчик частиц
\(N = 1\)

\textbf{Шаг 2. Запуск частицы:} - Вычислить \(R_{max}\) --- максимальный
радиус кластера - Запустить с окружности:
\(R_{start} = R_{max} + \Delta\) - Случайный угол
\(\alpha \in [0, 2\pi)\)

\textbf{Шаг 3. Случайное блуждание:} - Выбрать случайное направление (4
варианта) - Сделать шаг, проверить границы сетки

\textbf{Шаг 4. Проверка:} - Контакт с кластером --- прилипание - Уход
дальше \(2R_{max}\) --- уничтожение частицы
\end{frame}

\begin{frame}{4.2 Оптимизация вычислений}
\phantomsection\label{ux43eux43fux442ux438ux43cux438ux437ux430ux446ux438ux44f-ux432ux44bux447ux438ux441ux43bux435ux43dux438ux439}
\begin{itemize}
\tightlist
\item
  \textbf{Стартовая окружность:} запуск частиц рядом с кластером
\item
  \textbf{Зона уничтожения:} отсев частиц, ушедших слишком далеко
\item
  \textbf{Результат:} ускорение в десятки раз
\end{itemize}
\end{frame}

\section{5. Модификации
алгоритма}\label{ux43cux43eux434ux438ux444ux438ux43aux430ux446ux438ux438-ux430ux43bux433ux43eux440ux438ux442ux43cux430}

\begin{frame}{5.1 Вероятность прилипания}
\phantomsection\label{ux432ux435ux440ux43eux44fux442ux43dux43eux441ux442ux44c-ux43fux440ux438ux43bux438ux43fux430ux43dux438ux44f}
\textbf{Химически-ограниченная агрегация:}

\begin{longtable}[]{@{}ll@{}}
\toprule\noalign{}
Число соседей \(k\) & Вероятность \(p_k\) \\
\midrule\noalign{}
\endhead
1 & 0.01 \\
2 & 0.10 \\
3 & 0.90 \\
4 & 1.00 \\
\bottomrule\noalign{}
\end{longtable}

При отказе --- частица отступает и продолжает блуждание
\end{frame}

\begin{frame}{5.2 Бессеточная модель}
\phantomsection\label{ux431ux435ux441ux441ux435ux442ux43eux447ux43dux430ux44f-ux43cux43eux434ux435ux43bux44c}
\begin{itemize}
\tightlist
\item
  Координаты --- вещественные числа \(x, y \in \mathbb{R}\)
\item
  Шаг блуждания: случайный угол \(\alpha\), фиксированная длина \(h\)
\item
  Прилипание: расстояние до кластера \(\le 2r\)
\item
  Дает более изотропные кластеры
\end{itemize}
\end{frame}

\section{6. Другие модели
агрегации}\label{ux434ux440ux443ux433ux438ux435-ux43cux43eux434ux435ux43bux438-ux430ux433ux440ux435ux433ux430ux446ux438ux438}

\begin{frame}{6.1 Баллистическая агрегация}
\phantomsection\label{ux431ux430ux43bux43bux438ux441ux442ux438ux447ux435ux441ux43aux430ux44f-ux430ux433ux440ux435ux433ux430ux446ux438ux44f}
\begin{itemize}
\tightlist
\item
  Частицы падают сверху по прямой
\item
  Случайные боковые смещения с вероятностью около \(0.2\)
\item
  Формируется плотный агрегат
\item
  Граница кластера --- фрактал
\end{itemize}
\end{frame}

\begin{frame}{6.2 Кластер-кластерная агрегация}
\phantomsection\label{ux43aux43bux430ux441ux442ux435ux440-ux43aux43bux430ux441ux442ux435ux440ux43dux430ux44f-ux430ux433ux440ux435ux433ux430ux446ux438ux44f}
\begin{itemize}
\tightlist
\item
  Изначально много одиночных частиц
\item
  Все кластеры движутся и слипаются
\item
  Коэффициент диффузии: \(D(M) = D_0 M^{-\gamma}\)
\item
  Самые рыхлые структуры (\(D \approx 1.4\))
\end{itemize}
\end{frame}

\section{7. Результаты}\label{ux440ux435ux437ux443ux43bux44cux442ux430ux442ux44b}

\begin{frame}{7.1 Сравнение размерностей}
\phantomsection\label{ux441ux440ux430ux432ux43dux435ux43dux438ux435-ux440ux430ux437ux43cux435ux440ux43dux43eux441ux442ux435ux439}
\begin{longtable}[]{@{}ll@{}}
\toprule\noalign{}
Модель & Фрактальная размерность \(D\) \\
\midrule\noalign{}
\endhead
DLA классическая & \(\approx 1.71\) \\
Химически-ограниченная & \(1.75\) --- \(1.85\) \\
Бессеточная DLA & \(\approx 1.68\) --- \(1.72\) \\
Кластер-кластерная & \(\approx 1.40\) --- \(1.50\) \\
\bottomrule\noalign{}
\end{longtable}
\end{frame}

\begin{frame}{7.2 Выводы по этапу}
\phantomsection\label{ux432ux44bux432ux43eux434ux44b-ux43fux43e-ux44dux442ux430ux43fux443}
\begin{itemize}
\tightlist
\item
  Разработаны алгоритмы для 6 моделей агрегации
\item
  Реализованы два метода расчета фрактальной размерности
\item
  Применены оптимизации для ускорения вычислений
\item
  Алгоритмическая база готова к программной реализации
\end{itemize}
\end{frame}

\section{8. Итоги}\label{ux438ux442ux43eux433ux438}

\begin{frame}{8.1 Заключение}
\phantomsection\label{ux437ux430ux43aux43bux44eux447ux435ux43dux438ux435}
\textbf{Что сделано:}

\begin{itemize}
\tightlist
\item
  Формализован математический аппарат
\item
  Разработаны пошаговые алгоритмы для всех моделей
\item
  Обоснован выбор вычислительных подходов
\end{itemize}

\textbf{Что дальше:}

\begin{itemize}
\tightlist
\item
  Программная реализация на Python
\item
  Визуализация кластеров
\item
  Сравнительный анализ результатов
\end{itemize}

\textbf{Главный результат:} эффективные алгоритмы, готовые к кодированию
\end{frame}




\end{document}
